% Generated draft: requires mathematical and source review.
\documentclass[11pt]{article}
\usepackage[margin=25mm]{geometry}
\usepackage{fontspec}
\setmainfont{Latin Modern Roman}
\usepackage{amsmath,amssymb,mathtools,unicode-math}
\setmathfont{Latin Modern Math}
\usepackage{enumitem,needspace,xurl,hyperref}
\setlength{\parindent}{0pt}
\setlength{\parskip}{5pt}
\setlength{\emergencystretch}{3em}
\newcommand{\N}{\mathbb{N}}
\newcommand{\Z}{\mathbb{Z}}
\newcommand{\Q}{\mathbb{Q}}
\newcommand{\R}{\mathbb{R}}
\newcommand{\C}{\mathbb{C}}
\newcommand{\F}{\mathbb{F}}
\newcommand{\A}{\mathbb{A}}
\newcommand{\PP}{\mathbb{P}}
\newcommand{\E}{\mathbb{E}}
\newcommand{\eps}{\varepsilon}
\newcommand{\dd}{\,\mathrm{d}}
\begin{document}
\section*{UnsolvedMath ID: 700893}
\textbf{Source problem code:} catalogue entry
\subsection{Definitions and mathematical statement}
For a group \(G\), let
\[
\operatorname{NSub}(G)=\{H\leq G: gHg^{-1}=H\text{ for every }g\in G\}.
\]
Its normal subgroup lattice has operations
\[
H\wedge K=H\cap K,\qquad H\vee K=\langle H\cup K\rangle.
\]
The literal reading of the supplied class is
\[
\mathcal N_{\mathrm{full}}=\{L:L\cong\operatorname{NSub}(G)\text{ as lattices for some group }G\}.
\]
Here groups are arbitrary; no finiteness assumption is supplied.

In the ordinary lattice language \((\wedge,\vee)\), a variety is a class axiomatized by identities between lattice terms. Thus the literal formulation of the question is: does there exist a set \(\Sigma\) of lattice identities such that, for every lattice \(L\),
\[
L\in\mathcal N_{\mathrm{full}}\quad\Longleftrightarrow\quad L\models\Sigma?
\]
All variables in each identity are universally quantified.

\textbf{Unresolved qualification.} The source also reports that its class \(N\) is an infinitely based quasivariety. A quasivariety is axiomatized by universally quantified implications of the form
\[
\left(\bigwedge_{i=1}^{m}s_i=t_i\right)\Longrightarrow s=t,
\]
where the expressions are lattice terms; being infinitely based means that no finite set of such axioms suffices. This is a supplied literature assertion, not a result verified here. The intended meaning of \(N\) requires checking: the class of full normal subgroup lattices and the class of lattices isomorphic to sublattices of normal subgroup lattices are distinct possible representability targets and must not be silently identified. The question displayed above records the literal full-lattice reading only. The supplied record also does not explicitly specify whether least and greatest elements are included as constants in the lattice language; the displayed formulation uses the ordinary language without those constants.
\subsection{Short English statement}
The supplied problem asks whether normal subgroup lattices of groups form a class definable by lattice identities. Its accompanying assertion about an infinitely based quasivariety requires clarification of whether full normal subgroup lattices or lattices embeddable into them are intended. The registry labels the problem open; this draft does not independently verify that status.
\subsection{Sources}
\begin{itemize}
\item Local UnsolvedMath registry, permanent ID \(700893\), source problem code \texttt{catalogue entry}, title \emph{Kourovka Notebook Problem 12.75}; supplied upstream revision \texttt{d371ee917ac9647fbb06cd6fc529deec6f3dda82}. The supplied registry status is \texttt{open}.
\item \emph{The Kourovka Notebook (Unsolved Problems in Group Theory)}, problem 12.75, 12th issue (1992); the supplied background also identifies the 21st issue edition (2026), \texttt{arXiv:1401.0300v48}. The collection description names E. I. Khukhro and V. D. Mazurov as editors.
\item Attribution is retained as supplied: the registry and background name D. M. Smirnov as proposer, with proposed year 1992; the statement itself begins with the attribution ``B. Jonsson.'' Their respective roles are not further explained in the supplied record.
\item C. Herrman and W. Poguntke, \emph{Algebra Univers.}, volume 4, no. 3 (1974), pp. 280--286, cited in the supplied statement for the assertion that \(N\) is an infinitely based quasivariety. Author spelling is retained as supplied; no article title or direct link is supplied. This citation and its precise scope have not been independently verified.
\end{itemize}
\end{document}
